% ============================================================================
% MATEMATICAS 2 - TEMA 1 - SEGUNDO PILOTO (ENFOQUE ACADEMICO)
% Exposicion para clase teorica: rigor, transiciones, ejemplos progresivos.
% NO sustituye temas/tema1.tex ni temas/tema1_piloto_didactico.tex.
% Compilar desde la raiz: latexmk -pdf temas/tema1_piloto_academico.tex
% Hereda la plantilla azul, 16:9 y 12pt de main.tex/config.tex.
% ============================================================================
\documentclass[main.tex]{subfiles}

\begin{document}

% Imagen nueva o no localizada: cajetin deliberadamente discreto.
\newcommand{\imagenpendiente}[3]{%
  \begin{center}
  \fcolorbox{mateBlueDark}{mateBlueInner}{%
    \begin{minipage}[c][#2][c]{#1}
      \centering\color{mateBlueDark}
      {\sffamily\bfseries Ilustración pendiente}\par\medskip
      {\sffamily\footnotesize #3}
    \end{minipage}}%
  \end{center}%
}

% Aprovechar originales del tema 1; si faltan, mostrar una indicacion.
% La doble ruta facilita la compilacion desde la raiz o desde temas/.
\newcommand{\imagenoriginal}[4]{%
  \IfFileExists{imagenes/tema1/#1}{%
    \includegraphics[width=#2,height=#3,keepaspectratio]{imagenes/tema1/#1}%
  }{%
    \IfFileExists{../imagenes/tema1/#1}{%
      \includegraphics[width=#2,height=#3,keepaspectratio]{../imagenes/tema1/#1}%
    }{%
      \imagenpendiente{#2}{#3}{#4}%
    }%
  }%
}

\portadatema{1}{Funciones de varias variables}

% ---- INTRODUCCION ---------------------------------------------------------
% 1
\begin{frame}{Índice del tema}
\begin{enumerate}
  \item Introducción: funciones de una y varias variables.
  \item Conceptos fundamentales: dominio, codominio e imagen.
  \item Determinación y representación de dominios.
  \item Representación gráfica de funciones de dos variables.
  \item Curvas de nivel: cálculo e interpretación.
  \item Funciones de tres o más variables.
  \item Ejercicios propuestos.
\end{enumerate}
\end{frame}

% 2
\begin{frame}{De una variable a varias variables}
En una función de una variable, un dato de entrada determina un valor:
\[
f:D\subseteq\mathbb R\longrightarrow\mathbb R,\qquad y=f(x).
\]
\medskip

En numerosos problemas de ingeniería el resultado depende de \textbf{varias magnitudes independientes}.
\medskip
\begin{itemize}
  \item Temperatura sobre una superficie: $T(x,y)$.
  \item Volumen de una caja: $V(x,y,z)=xyz$.
  \item Intensidad de una imagen: $I(i,j)$, con $i,j$ índices de píxeles.
\end{itemize}
\end{frame}

% 3
\begin{frame}{Definición: función real de varias variables}
Una \textbf{función real de $n$ variables} es una aplicación
\[
f:D\subseteq\mathbb R^n\longrightarrow\mathbb R
\]
que asigna a cada elemento $(x_1,\ldots,x_n)\in D$ un único valor
\[
u=f(x_1,\ldots,x_n).
\]
\medskip

Nos centraremos primero en el caso $n=2$, porque sus funciones pueden interpretarse geométricamente mediante superficies.
\end{frame}

% 4
\begin{frame}{Evaluación de una función de dos variables}
Consideremos
\[
f(x,y)=x^2+2y.
\]
Calcular el valor de la función consiste en sustituir \textbf{ambas} variables:
\[
f(1,2)=1^2+2\cdot2=5,\qquad
f(-2,1)=(-2)^2+2=6.
\]
\medskip

Distintas entradas pueden producir una misma salida; lo que exige la definición es que \textbf{cada entrada tenga una sola salida}.
\end{frame}

% ---- CONCEPTOS FUNDAMENTALES ---------------------------------------------
% 5
\begin{frame}{Dominio, codominio e imagen}
Sea $f:D\subseteq\mathbb R^2\to\mathbb R$.
\begin{itemize}
  \item \textbf{Dominio $D$:} pares $(x,y)$ para los que está definida.
  \item \textbf{Codominio:} conjunto de llegada fijado al definir $f$; aquí, $\mathbb R$.
  \item \textbf{Imagen:} valores efectivamente alcanzados:
\end{itemize}
\[
\operatorname{Im}(f)=\{f(x,y):(x,y)\in D\}.
\]
\medskip

En general, $\operatorname{Im}(f)$ es un subconjunto del codominio.
\end{frame}

% 6
\begin{frame}{Ejemplo: el codominio puede ser mayor que la imagen}
Consideremos $f:\mathbb R^2\to\mathbb R$ dada por
\[
f(x,y)=x^2+y^2.
\]
\begin{itemize}
  \item Dominio: $D=\mathbb R^2$.
  \item Codominio: $\mathbb R$.
  \item Imagen: $[0,+\infty)$.
\end{itemize}
\medskip

La función nunca toma valores negativos y alcanza el mínimo $0$ en $(0,0)$.
\end{frame}

% ---- DOMINIOS -------------------------------------------------------------
% 7
\begin{frame}{Determinación del dominio natural}
Cuando una función se presenta únicamente mediante su expresión, buscamos el mayor conjunto de entradas donde produce valores reales.
\medskip

Restricciones habituales:
\begin{itemize}
  \item Denominadores: no pueden anularse.
  \item Raíces de índice par: el radicando debe ser $\geq0$.
  \item Logaritmos: el argumento debe ser $>0$.
  \item $\arcsin(t)$ y $\arccos(t)$: $-1\leq t\leq1$.
\end{itemize}
\medskip

Cuando aparecen varias condiciones, el dominio es su \textbf{intersección}.
\end{frame}

% 8
\begin{frame}{Ejemplo 1: dominio de una función polinómica}
Sea
\[
f(x,y)=x^2+3xy-y^3.
\]
Las operaciones de suma, producto y potencias enteras no negativas están definidas para cualquier par real.
\[
\boxed{D=\mathbb R^2.}
\]
\medskip

Este es el caso de partida: \textbf{no hay restricciones adicionales}.
\end{frame}

% 9
\begin{frame}{Ejemplo 2: una restricción de tipo radical}
Sea
\[
f(x,y)=\sqrt{9-x^2-y^2}.
\]
Para que el resultado sea real:
\[
9-x^2-y^2\geq0
\quad\Longleftrightarrow\quad x^2+y^2\leq9.
\]
\[
\boxed{D=\{(x,y)\in\mathbb R^2:x^2+y^2\leq9\}.}
\]
El dominio es un disco cerrado de radio $3$.
\end{frame}

% 10
\begin{frame}{Interpretación geométrica del dominio}
\begin{columns}[c]
  \begin{column}{.49\textwidth}
    \centering
    \begin{tikzpicture}[scale=.68]
      \fill[mateBlueInner] (0,0) circle (3);
      \draw[->,gray] (-3.6,0)--(3.75,0) node[right]{$x$};
      \draw[->,gray] (0,-3.6)--(0,3.7) node[above]{$y$};
      \draw[very thick,mateBlueDark] (0,0) circle (3);
      \node[below right] at (3,0) {$3$};
      \node[below left] at (-3,0) {$-3$};
    \end{tikzpicture}
  \end{column}
  \begin{column}{.49\textwidth}
    \textbf{Dominio de}
    \[
      \sqrt{9-x^2-y^2}.
    \]
    La frontera $x^2+y^2=9$ \textbf{está incluida}, porque la raíz de cero existe.
    \medskip

    Fuera del disco, el radicando es negativo.
  \end{column}
\end{columns}
\end{frame}

% 11
\begin{frame}{Ejemplo 3: dominio de una función racional}
Sea
\[
f(x,y)=\frac{3x-5y}{y-x^2}.
\]
El denominador no puede ser cero:
\[
y-x^2\ne0 \quad\Longleftrightarrow\quad y\ne x^2.
\]
Por tanto,
\[
\boxed{D=\{(x,y)\in\mathbb R^2:y\ne x^2\}.}
\]
\medskip

Geométricamente, es todo el plano salvo los puntos de la parábola $y=x^2$.
\end{frame}

% 12
\begin{frame}{Ejemplo 4: dominio de una función logarítmica}
Sea
\[
f(x,y)=\ln(y-x^2).
\]
Es preciso que el argumento sea estrictamente positivo:
\[
y-x^2>0 \quad\Longleftrightarrow\quad y>x^2.
\]
\[
\boxed{D=\{(x,y)\in\mathbb R^2:y>x^2\}.}
\]
\medskip

Es la región situada \textbf{por encima} de la parábola; la propia parábola queda excluida.
\end{frame}

% 13
\begin{frame}{Ejemplo 5: restricciones simultáneas}
Consideremos
\[
f(x,y)=\frac{\sqrt{9-x^2-y^2}}{\ln(y+2)}.
\]
Deben verificarse \textbf{tres condiciones}:
\begin{align*}
9-x^2-y^2&\geq0 &&\text{(radical)},\\
y+2&>0 &&\text{(logaritmo)},\\
\ln(y+2)&\ne0 &&\text{(denominador)}.
\end{align*}
\end{frame}

% 14
\begin{frame}{Ejemplo 5: intersección de las restricciones}
Las condiciones anteriores equivalen a
\[
x^2+y^2\leq9,\qquad y>-2,\qquad y\ne-1.
\]
Por tanto,
\[
\boxed{D=\left\{(x,y)\in\mathbb R^2:
\begin{array}{l}
x^2+y^2\leq9,\\
y>-2,\quad y\ne-1
\end{array}\right\}.}
\]
\medskip

\textbf{Interpretación:} la parte del disco que está por encima de $y=-2$, excluyendo los puntos de la recta $y=-1$.
\end{frame}

% 15
\begin{frame}{Síntesis: cómo determinar un dominio}
\begin{enumerate}
  \item Identificar las operaciones que imponen restricciones.
  \medskip
  \item Expresar cada restricción mediante igualdades o desigualdades.
  \medskip
  \item Calcular la intersección de las regiones válidas.
  \medskip
  \item Representar geométricamente el resultado cuando sea posible.
\end{enumerate}
\medskip

Las restricciones se deducen de la expresión, pero la representación geométrica ayuda a \textbf{comprobar} su interpretación.
\end{frame}

% ---- REPRESENTACION GRAFICA -----------------------------------------------
% 16
\begin{frame}{Del dominio a la representación gráfica}
Conocer el dominio no basta para describir cómo varía una función.
\medskip

Para $f:D\subseteq\mathbb R^2\to\mathbb R$, su \textbf{gráfica} es
\[
\boxed{G_f=\{(x,y,f(x,y)):(x,y)\in D\}.}
\]
\medskip
\begin{itemize}
  \item El dominio se representa en $\mathbb R^2$.
  \item La gráfica está contenida en $\mathbb R^3$.
\end{itemize}
\medskip

En muchas funciones de dos variables, la gráfica es una superficie.
\end{frame}

% 17
\begin{frame}{Ejemplo 1: una superficie plana}
Sea $f(x,y)=x+2y$.
\[
z=x+2y.
\]
Su gráfica es un plano que contiene, entre otros, los puntos
\[
(0,0,0),\qquad (1,0,1),\qquad (0,1,2).
\]
\medskip

\begin{center}
  \imagenpendiente{.65\linewidth}{27mm}{Plano $z=x+2y$, con ejes y tres puntos señalados.}
\end{center}
\end{frame}

% 18
\begin{frame}{Ejemplo 2: una superficie curva}
Sea
\[
f(x,y)=x+y^2,\qquad D=\mathbb R^2.
\]
La gráfica es la superficie $z=x+y^2$. Su altura crece linealmente con $x$ y de forma cuadrática con $y$.
\begin{center}
\imagenoriginal{s9_2.png}{.82\linewidth}{39mm}{Superficie original de $z=x+y^2$.}
\end{center}
\end{frame}

% 19
\begin{frame}{Ejemplo 3: la semiesfera superior}
La función
\[
f(x,y)=\sqrt{9-x^2-y^2}
\]
satisface
\[
x^2+y^2+z^2=9,\qquad z\geq0.
\]
Su gráfica es la \textbf{semiesfera superior} de radio $3$.
\medskip

Su dominio es un disco en el plano; su gráfica es una superficie en el espacio.
\end{frame}

% ---- CURVAS DE NIVEL ------------------------------------------------------
% 20
\begin{frame}{Por qué estudiamos las curvas de nivel}
La gráfica de una función de dos variables requiere una representación tridimensional.
\medskip

En cartografía, visualización científica y análisis de funciones de coste interesa a menudo representar esa información \textbf{en un plano}.
\medskip

La idea es agrupar los puntos del dominio en los que $f$ toma \textbf{el mismo valor}.
\end{frame}

% 21
\begin{frame}{Definición de curva de nivel}
Sea $f:D\subseteq\mathbb R^2\to\mathbb R$ y sea $k\in\mathbb R$.
\medskip

El \textbf{conjunto de nivel} correspondiente a $k$ es
\[
\boxed{C_k=\{(x,y)\in D:f(x,y)=k\}.}
\]
\medskip

Habitualmente, estos conjuntos son curvas; también pueden reducirse a un punto, ser vacíos o presentar otras geometrías.
\end{frame}

% 22
\begin{frame}{Ejemplo 1: curvas de nivel rectilíneas}
Para $f(x,y)=x+y$, el nivel $k$ verifica
\[
x+y=k\quad\Longleftrightarrow\quad y=k-x.
\]
\medskip

Por tanto, las curvas de nivel son \textbf{rectas paralelas}.
\[
\begin{array}{c|c}
k&\text{curva de nivel}\\\hline
-1&y=-1-x\\
0&y=-x\\
1&y=1-x
\end{array}
\]
\end{frame}

% 23
\begin{frame}{Ejemplo 2: curvas de nivel circulares}
Para $f(x,y)=x^2+y^2$, el nivel $k$ verifica
\[
x^2+y^2=k.
\]
\begin{itemize}
  \item Si $k>0$: circunferencia de radio $\sqrt{k}$.
  \item Si $k=0$: únicamente el punto $(0,0)$.
  \item Si $k<0$: conjunto vacío.
\end{itemize}
\medskip

La ecuación permite estudiar \textbf{todos los niveles}, incluso aquellos que no producen una curva.
\end{frame}

% 24
\begin{frame}{Representación de varios niveles}
\begin{columns}[c]
  \begin{column}{.48\textwidth}
    \centering
    \begin{tikzpicture}[scale=.73]
      \draw[->,gray] (-2.55,0)--(2.7,0) node[right] {$x$};
      \draw[->,gray] (0,-2.55)--(0,2.7) node[above] {$y$};
      \draw[very thick,mateBlueDark] (0,0) circle (2);
      \draw[very thick,mateBlueOuter] (0,0) circle (1);
      \fill[mateBlueDark] (0,0) circle (.075);
      \node[below right] at (2,0) {$k=4$};
      \node[above right] at (.78,.70) {$k=1$};
      \node[below left] at (0,0) {$k=0$};
    \end{tikzpicture}
  \end{column}
  \begin{column}{.5\textwidth}
    Para $f(x,y)=x^2+y^2$:
    \[
    C_1:\ x^2+y^2=1,
    \]
    \[
    C_4:\ x^2+y^2=4.
    \]
    Los niveles de mayor valor están más alejados del origen.
  \end{column}
\end{columns}
\end{frame}

% 25
\begin{frame}{Ejemplo 3: curvas de nivel hiperbólicas}
Sea
\[
f(x,y)=x^2-y^2.
\]
Sus niveles vienen dados por $x^2-y^2=k$.
\begin{itemize}
  \item $k>0$: hipérbolas abiertas hacia la dirección del eje $x$.
  \item $k<0$: hipérbolas abiertas hacia la dirección del eje $y$.
  \item $k=0$: dos rectas, $y=x$ e $y=-x$.
\end{itemize}
\medskip

Los niveles permiten distinguir la geometría de esta función de la de $x^2+y^2$.
\end{frame}

% 26
\begin{frame}{Ejemplo 4: curvas de nivel de una semiesfera}
Para
\[
f(x,y)=\sqrt{100-x^2-y^2}
\]
el nivel $c$ satisface
\[
c=\sqrt{100-x^2-y^2}
\quad\Longrightarrow\quad
\boxed{x^2+y^2=100-c^2.}
\]
\medskip
\begin{itemize}
  \item $0\leq c<10$: circunferencia de radio $\sqrt{100-c^2}$.
  \item $c=10$: punto $(0,0)$.
  \item $c<0$ o $c>10$: conjunto vacío.
\end{itemize}
\end{frame}

% 27
\begin{frame}{Aplicación: funciones de coste en Informática}
Un problema de ajuste de parámetros puede conducir a una función
\[
L(w_1,w_2)=(w_1-2)^2+2(w_2+1)^2.
\]
\medskip

Sus curvas de nivel cumplen
\[
(w_1-2)^2+2(w_2+1)^2=k.
\]
Para $k>0$, son \textbf{elipses} centradas en $(2,-1)$.
\medskip

Esta representación se utiliza para estudiar visualmente funciones objetivo en optimización y aprendizaje automático.
\end{frame}

% 28
\begin{frame}{Aplicación: intensidad de una imagen digital}
Una imagen en escala de grises puede modelarse como
\[
I:\{0,\ldots,M-1\}\times\{0,\ldots,N-1\}
\longrightarrow\{0,\ldots,255\}.
\]
\medskip

La intensidad depende de dos índices $(i,j)$, pero, a diferencia de nuestros modelos continuos, su \textbf{dominio es discreto}.
\medskip

\imagenpendiente{.71\linewidth}{24mm}{Imagen de grises y matriz de intensidades de una zona ampliada.}
\end{frame}

% ---- GENERALIZACION -------------------------------------------------------
% 29
\begin{frame}{Funciones de tres o más variables}
La definición se extiende a cualquier número de variables:
\[
f:D\subseteq\mathbb R^n\longrightarrow\mathbb R.
\]
\medskip

Por ejemplo,
\[
V(x,y,z)=xyz
\]
asigna un volumen a tres dimensiones positivas de una caja.
\medskip

La gráfica de una función de tres variables estaría contenida en $\mathbb R^4$; por ello adquieren importancia las representaciones mediante \textbf{conjuntos de nivel}.
\end{frame}

% 30
\begin{frame}{Ejemplo: superficies de nivel}
Consideremos
\[
f(x,y,z)=x^2+y^2.
\]
Para el nivel $k=4$:
\[
x^2+y^2=4.
\]
\medskip

Es una \textbf{superficie cilíndrica} de radio $2$, paralela al eje $z$: la coordenada $z$ es libre.
\medskip

Para $k=0$, el conjunto de nivel es todo el eje $z$, no un único punto.
\end{frame}

% ---- EJEMPLO INTEGRADOR ---------------------------------------------------
% 31
\begin{frame}{Ejemplo integrador: dominio, imagen y niveles}
Sea $f(x,y)=\sqrt{4-x^2-y^2}$.
\begin{enumerate}
  \item \textbf{Dominio:}
  \[
  D=\{(x,y):x^2+y^2\leq4\}.
  \]
  \item \textbf{Imagen:} $[0,2]$.
  \item \textbf{Nivel $k=1$:}
  \[
  \sqrt{4-x^2-y^2}=1
  \quad\Longleftrightarrow\quad
  x^2+y^2=3.
  \]
\end{enumerate}
\end{frame}

% ---- EJERCICIOS PROPUESTOS ------------------------------------------------
% 32
\begin{frame}{Ejercicios básicos}
\textbf{Objetivo:} aplicar directamente los conceptos.
\medskip
\begin{enumerate}
  \item Para $f(x,y)=2x+y$, calcula $f(1,2)$ y $f(-1,3)$.
  \medskip
  \item Determina el dominio natural de
  \[
  g(x,y)=\sqrt{4-x^2-y^2}.
  \]
  \item Describe los niveles $k=0,1,4$ de
  \[
  q(x,y)=x^2+y^2.
  \]
\end{enumerate}
\end{frame}

% 33
\begin{frame}{Ejercicios intermedios}
\textbf{Objetivo:} resolver problemas de varios pasos.
\medskip
\begin{enumerate}
  \item Determina y representa el dominio de
  \[
  f(x,y)=\ln(y-x^2).
  \]
  \item Halla el dominio de
  \[
  g(x,y)=\frac{3x-5y}{y-x^2}.
  \]
  \item Clasifica los niveles $k=-1,0,1$ de $h(x,y)=x^2-y^2$.
\end{enumerate}
\end{frame}

% 34
\begin{frame}{Ejercicios avanzados}
\textbf{Objetivo:} integrar varias restricciones y generalizar.
\medskip
\begin{enumerate}
  \item Obtén el dominio de
  \[
  F(x,y,z)=\arcsin(xy)e^{3z}+e^{(y+z)\ln(x+z)}.
  \]
  \item Describe los conjuntos de nivel para $k=0,1,4$ de
  \[
  G(x,y,z)=x^2+y^2.
  \]
\end{enumerate}
\medskip

\textit{Ejercicios de ampliación respecto de la práctica básica del tema.}
\end{frame}

% 35
\begin{frame}{Síntesis y relación con los temas siguientes}
Hemos estudiado:
\begin{itemize}
  \item Cómo se definen las funciones reales de varias variables.
  \item Cómo determinar su dominio y su imagen.
  \item Cómo representar funciones de dos variables en el espacio.
  \item Cómo calcular e interpretar curvas y superficies de nivel.
\end{itemize}
\medskip

\textbf{Continuación del cálculo:} tras repasar la derivación en una variable, estudiaremos cómo describir la variación de estas funciones respecto de cada una de sus variables.
\end{frame}

\end{document}
